\chapter{2008年安徽卷（理）}

\section{单选题}

\begin{problem}
复数 \(\i^3 (1 + \i)^2 =\)（\quad）

\choices
    {\(2\)}
    {\(-2\)}
    {\(2\i\)}
    {\(-2\i\)}
\end{problem}

\begin{answer}
A
\end{answer}

\begin{solution}
因为 \((1+\i)^2=2\i\)，且 \(\i^3=-\i\)，所以
\[
\i^3(1+\i)^2=(-\i)(2\i)=2.
\]
故选 A.
\end{solution}

\begin{problem}
集合 \(A = \{y \in \R \mid y = \lg x, x > 1\}\)，\(B = \{-2, -1, 1, 2\}\) 则下列结论正确的是（\quad）

\choices
    {\(A \cap B = \{-2, -1\}\)}
    {\((\complement_{\R}A)\cup B = (-\infty ,0)\)}
    {\(A \cup B = (0, +\infty)\)}
    {\((\complement_{\R}A)\cap B = \{-2, -1\}\)}
\end{problem}

\begin{answer}
D
\end{answer}

\begin{solution}
当 \(x>1\) 时，\(\lg x>0\)，故 \(A=(0,+\infty)\). 因此
\[
A\cap B=\{1,2\},\qquad \complement_{\R}A=(-\infty,0],
\]
从而 \((\complement_{\R}A)\cap B=\{-2,-1\}\). 其余三个结论均不成立，故选 D.
\end{solution}

\begin{problem}
在平行四边形 \(ABCD\) 中，\(AC\) 为一条对角线，若 \(\overrightarrow{AB} = (2, 4)\)，\(\overrightarrow{AC} = (1, 3)\)，则 \(\overrightarrow{BD} =\)（\quad）

\choices
    {\((-2, - 4)\)}
    {\((-3, - 5)\)}
    {\((3, 5)\)}
    {\((2, 4)\)}
\end{problem}

\begin{answer}
B
\end{answer}

\begin{solution}
在平行四边形中，\(\overrightarrow{BC}=\overrightarrow{AC}-\overrightarrow{AB}\)，所以
\[
\overrightarrow{BD}=\overrightarrow{BC}-\overrightarrow{AB}
=\overrightarrow{AC}-2\overrightarrow{AB}
=(1,3)-2(2,4)=(-3,-5).
\]
故选 B.
\end{solution}

\begin{problem}
已知 \(m\)，\(n\) 是两条不同直线，\(\alpha\)，\(\beta\)，\(\gamma\) 是三个不同平面，下列命题中正确的是（\quad）

\choices
    {若 \(m \parallel \alpha, n \parallel \alpha\)，则 \(m \parallel n\)}
    {若 \(\alpha \perp \gamma\)，\(\beta \perp \gamma\)，则 \(\alpha \parallel \beta\)}
    {若 \(m \parallel \alpha\)，\(m \parallel \beta\)，则 \(\alpha \parallel \beta\)}
    {若 \(m \perp \alpha\)，\(n \perp \alpha\)，则 \(m \parallel n\)}
\end{problem}

\begin{answer}
D
\end{answer}

\begin{solution}
若两条直线都平行于同一平面，它们仍可能相交，故 \(A\) 不正确；两个都垂直于同一平面的平面不一定平行，故 \(B\) 不正确；一条直线同时平行于两个相交平面也是可能的，故 \(C\) 不正确.

直线垂直于平面，就平行于该平面的法线方向. 因而 \(m\) 与 \(n\) 都平行于平面 \(\alpha\) 的法线，故 \(m\parallel n\)，选 D.
\end{solution}

\begin{problem}
将函数 \( y = \sin \left(2x + \frac{\pi}{3}\right) \) 的图象按向量 \( \bs{a} \) 平移后所得的图象关于点 \( \left(-\frac{\pi}{12}, 0\right) \) 中心对称，则向量 \( \bs{a} \) 的坐标可能为（\quad）

\choices
    {\( \left(-\frac{\pi}{12},0\right) \)}
    {\( \left(-\frac{\pi}{6},0\right) \)}
    {\(\left(\frac{\pi}{12},0\right)\)}
    {\(\left(\frac{\pi}{6},0\right)\)}
\end{problem}

\begin{answer}
C
\end{answer}

\begin{solution}
设平移向量为 \(\bs{a}=(h,k)\)，平移后的函数为
\[
y=\sin\left(2(x-h)+\frac{\pi}{3}\right)+k.
\]
正弦曲线的中心对称点的纵坐标等于其竖直平移量，所以题给中心对称点的纵坐标为 \(0\) 时必须有 \(k=0\). 又因为该点在图象上，故存在 \(t\in\Z\)，使
\[
2\left(-\frac{\pi}{12}-h\right)+\frac{\pi}{3}=t\pi.
\]
于是 \(h=\frac{\pi}{12}-\frac{t\pi}{2}\). 取 \(t=0\)，得 \(\bs{a}=\left(\frac{\pi}{12},0\right)\)，故选 C.
\end{solution}

\begin{problem}
设 \((1 + x)^{8} = a_{0} + a_{1}x + \dots +a_{8}x^{8}\)，则 \(a_0\)，\(a_1\)，\(\dots\)，\(a_8\) 中奇数的个数为（\quad）

\choices
    {\(2\)}
    {\(3\)}
    {\(4\)}
    {\(5\)}
\end{problem}

\begin{answer}
A
\end{answer}

\begin{solution}
由二项式定理，\(a_k=\mathrm{C}_8^k\). 逐项计算得
\[
(a_0,a_1,\ldots,a_8)=(1,8,28,56,70,56,28,8,1).
\]
其中只有 \(a_0\) 和 \(a_8\) 是奇数，共有 \(2\) 个，故选 A.
\end{solution}

\begin{problem}
\(a < 0\) 是方程 \(ax^2 + 2x + 1 = 0\) 至少有一个负数根的（\quad）

\choices
    {必要不充分条件}
    {充分不必要条件}
    {充分必要条件}
    {既不充分也不必要条件}
\end{problem}

\begin{answer}
B
\end{answer}

\begin{solution}
若 \(a<0\)，则方程的判别式为 \(\Delta=4-4a>0\)，且两根之积为 \(\frac{1}{a}<0\)，所以方程有一个负根，故 \(a<0\) 是充分条件.

但它不是必要条件，例如取 \(a=1\)，方程为 \(x^2+2x+1=0\)，根为 \(-1\)，仍有负根. 因此选 B，即充分不必要条件.
\end{solution}

\begin{problem}
若过点 \(A(4,0)\) 的直线 \(l\) 与曲线 \((x - 2)^2 + y^2 = 1\) 有公共点，则直线 \(l\) 的斜率的取值范围为（\quad）

\choices
    {\([- \sqrt{3}, \sqrt{3} ]\)}
    {\((- \sqrt{3}, \sqrt{3} ]\)}
    {\(\left[-\frac{\sqrt{3}}{3}, \frac{\sqrt{3}}{3}\right]\)}
    {\(\left(-\frac{\sqrt{3}}{3}, \frac{\sqrt{3}}{3}\right)\)}
\end{problem}

\begin{answer}
C
\end{answer}

\begin{solution}
设直线 \(l\) 的斜率为 \(k\)，则其方程为 \(y=k(x-4)\)，即
\[
kx-y-4k=0.
\]
圆心 \((2,0)\) 到直线的距离不超过半径 \(1\)，所以
\[
\frac{|2k-4k|}{\sqrt{k^2+1}}\leqslant1.
\]
化简得 \(4k^2\leqslant k^2+1\)，即 \(|k|\leqslant\frac{\sqrt{3}}{3}\). 因此斜率的取值范围为 \(\left[-\frac{\sqrt{3}}{3},\frac{\sqrt{3}}{3}\right]\)，故选 C.
\end{solution}

\begin{problem}
在同一平面直角坐标系中，函数 \( y = g(x) \) 的图象与 \( y = \e^{x} \) 的图象关于直线 \( y = x \) 对称. 而函数 \( y = f(x) \) 的图象与 \( y = g(x) \) 的图象关于 \( y \) 轴对称. 若 \( f(m) = -1 \)，则 \( m \) 的值是（\quad）

\choices
    {\(-\e\)}
    {\(-\frac{1}{\e}\)}
    {\(\e\)}
    {\(\frac{1}{\e}\)}
\end{problem}

\begin{answer}
B
\end{answer}

\begin{solution}
图象关于直线 \(y=x\) 对称意味着互为反函数，因此 \(g(x)=\ln x\)（\(x>0\)）. 图象关于 \(y\) 轴对称意味着 \(f(x)=g(-x)=\ln(-x)\)（\(x<0\)）.

由 \(f(m)=-1\)，得 \(\ln(-m)=-1\)，所以 \(-m=\e^{-1}\)，即 \(m=-\frac{1}{\e}\)，故选 B.
\end{solution}

\begin{problem}
设两个正态分布 \(N(\mu_1, \sigma_1^2)\)（\(\sigma_1 > 0\)）和 \(N(\mu_2, \sigma_2^2)\)（\(\sigma_2 > 0\)）的密度函数图象如图所示，则有（\quad）

\bitmapfigure[max width=0.42\linewidth]{img/2008/anhui_science/q10_fig1.png}

\choices
    {\(\mu_{1} < \mu_{2}, \sigma_{1} < \sigma_{2}\)}
    {\(\mu_{1} < \mu_{2}\)，\(\sigma_{1} > \sigma_{2}\)}
    {\(\mu_{1} > \mu_{2}\)，\(\sigma_{1} < \sigma_{2}\)}
    {\(\mu_{1} > \mu_{2}\)，\(\sigma_{1} > \sigma_{2}\)}
\end{problem}

\begin{answer}
A
\end{answer}

\begin{solution}
正态分布密度曲线的对称轴为 \(x=\mu\)，曲线在对称轴处取得最大值 \(\frac{1}{\sqrt{2\pi}\sigma}\). 图中第一条曲线的对称轴在左侧，且峰值较高；第二条曲线的对称轴在右侧，且峰值较低.

因此 \(\mu_1<\mu_2\)，并且 \(\frac{1}{\sigma_1}>\frac{1}{\sigma_2}\)，即 \(\sigma_1<\sigma_2\)，故选 A.
\end{solution}

\begin{problem}
若函数 \(f(x)\)，\(g(x)\) 分别是 \( \R \) 上的奇函数，偶函数，且满足 \( f(x) - g(x) = \e^{x} \)，则有（\quad）

\choices
    {\(f(2) < f(3) < g(0)\)}
    {\(g(0) < f(3) < f(2)\)}
    {\(f(2) < g(0) < f(3)\)}
    {\(g(0) < f(2) < f(3)\)}
\end{problem}

\begin{answer}
D
\end{answer}

\begin{solution}
将条件中的 \(x\) 换成 \(-x\)，由 \(f\) 为奇函数、\(g\) 为偶函数，得
\[
-f(x)-g(x)=\e^{-x}.
\]
与原式 \(f(x)-g(x)=\e^x\) 相加、相减，得到
\[
f(x)=\frac{\e^x-\e^{-x}}{2},\qquad
g(x)=-\frac{\e^x+\e^{-x}}{2}.
\]
所以 \(g(0)=-1\)，且 \(f(2)>0\)，\(f(3)>f(2)\)，从而 \(g(0)<f(2)<f(3)\)，故选 D.
\end{solution}

\begin{problem}
\(12\) 名同学合影，站成前排 \(4\) 人后排 \(8\) 人，现摄影师要从后排 \(8\) 人中抽 \(2\) 人调整到前排，若其他人的相对顺序不变，则不同调整方法的总数是（\quad）

\choices
    {\(\mathrm{C}_8^2\mathrm{A}_3^2\)}
    {\(\mathrm{C}_8^2\mathrm{A}_6^6\)}
    {\(\mathrm{C}_8^2\mathrm{A}_6^2\)}
    {\(\mathrm{C}_8^2\mathrm{A}_5^2\)}
\end{problem}

\begin{answer}
D
\end{answer}

\begin{solution}
先从后排 \(8\) 人中选出调整到前排的 \(2\) 人，有 \(\mathrm{C}_8^2\) 种方法. 原前排的 \(4\) 人相对顺序不变，在这 \(4\) 人形成的 \(5\) 个位置中选出两个位置，并安排选出的两人，有 \(\mathrm{A}_5^2\) 种方法.

因此不同调整方法的总数为 \(\mathrm{C}_8^2\mathrm{A}_5^2\)，故选 D.
\end{solution}

\section{填空题}

\begin{problem}
函数 \( f(x) = \frac{\sqrt{|x - 2| - 1}}{\log_2(x - 1)} \) 的定义域为\fillinblank{}.
\end{problem}

\begin{answer}
\( [3, +\infty) \)
\end{answer}

\begin{solution}
根式有意义要求 \(|x-2|-1\geqslant0\)，即 \(x\leqslant1\) 或 \(x\geqslant3\). 同时对数要求 \(x-1>0\)，且分母不能为零，故 \(x>1\) 且 \(x-1\neq1\). 取交集得函数的定义域为 \([3,+\infty)\).
\end{solution}

\begin{problem}
在数列 \(\{a_{n}\}\) 中，\(a_{n} = 4n - \frac{5}{2}\)，\(a_{1} + a_{2} + \cdots + a_{n} = an^{2} + bn\)，\(n \in \N^*\)，其中 \(a\)，\(b\) 为常数，则 \(\displaystyle\lim_{n \to \infty} \frac{a^{n} - b^{n}}{a^{n} + b^{n}}\) 的值是\fillinblank{}.
\end{problem}

\begin{answer}
\(1\)
\end{answer}

\begin{solution}
数列前 \(n\) 项和为
\[
\sum_{k=1}^{n}\left(4k-\frac52\right)=2n(n+1)-\frac52n=2n^2-\frac12n.
\]
所以 \(a=2\)，\(b=-\frac12\). 因而
\[
\lim_{n\to\infty}\frac{a^n-b^n}{a^n+b^n}
=\lim_{n\to\infty}\frac{1-\left(-\frac14\right)^n}{1+\left(-\frac14\right)^n}=1.
\]
\end{solution}

\begin{problem}
若 \(A\) 为不等式组 \(\begin{cases}
x \leqslant 0,\\
y \geqslant 0,\\
y - x \leqslant 2
\end{cases}\) 表示的平面区域，则当 \(a\) 从 \(-2\) 连续变化到 1 时，动直线 \(x + y = a\) 扫过 \(A\) 中的那部分区域的面积为\fillinblank{}.
\end{problem}

\begin{answer}
\(\frac{7}{4}\)
\end{answer}

\begin{solution}
由 \(x\leqslant0\)、\(y\geqslant0\) 和 \(y\leqslant x+2\)，可知区域 \(A\) 是顶点为 \((-2,0)\)、\((0,0)\)、\((0,2)\) 的三角形，面积为 \(2\).

当 \(-2\leqslant a\leqslant1\) 时，直线 \(x+y=a\) 从顶点 \((-2,0)\) 扫到直线 \(x+y=1\). 未被扫过的部分是由 \((0,1)\)、\((0,2)\)、\(\left(-\frac12,\frac32\right)\) 围成的三角形，面积为 \(\frac12\times1\times\frac12=\frac14\). 所以扫过区域的面积为
\[
2-\frac14=\frac74.
\]
\end{solution}

\begin{problem}
已知 \(A\)，\(B\)，\(C\)，\(D\) 在同一个球面上，\(AB \perp\) 平面 \(BCD\)，\(BC \perp CD\)，若 \(AB = 6\)，\(AC = 2\sqrt{13}\)，\(AD = 8\)，则 \(B\)，\(C\) 两点间的球面距离是\fillinblank{}.
\end{problem}

\begin{answer}
\(\frac{4}{3}\pi\)
\end{answer}

\begin{solution}
因为 \(AB\perp\) 平面 \(BCD\)，所以
\[
BC=\sqrt{AC^2-AB^2}=\sqrt{52-36}=4,\qquad
BD=\sqrt{AD^2-AB^2}=\sqrt{64-36}=2\sqrt7.
\]
又因为 \(BC\perp CD\)，所以在平面 \(BCD\) 内，\(\angle BCD=90^{\circ}\)，从而 \(BD\) 是该平面内截得圆的直径，截圆半径为 \(\frac{BD}{2}=\sqrt7\).

球心到平面 \(BCD\) 的距离为 \(\frac{AB}{2}=3\)，故球半径满足
\[
R^2=\left(\sqrt7\right)^2+3^2=16,
\]
即 \(R=4\). 设 \(B,C\) 两点所对的较小圆心角为 \(\theta\)，则
\[
BC=2R\sin\frac{\theta}{2}
\quad\Longrightarrow\quad
4=8\sin\frac{\theta}{2},
\]
所以 \(\theta=\frac\pi3\). 因而两点间的球面距离为
\[
R\theta=4\cdot\frac\pi3=\frac{4\pi}{3}.
\]
\end{solution}

\section{解答题}

\begin{problem}
\begin{enumerate}
\item 已知函数 \( f(x) = \cos \left(2x - \frac{\pi}{3}\right) + 2\sin \left(x - \frac{\pi}{4}\right)\sin \left(x + \frac{\pi}{4}\right) \). 求函数 \( f(x) \) 的最小正周期和图象的对称轴方程；
\item 求函数 \( f(x) \) 在区间 \( \left[-\frac{\pi}{12}, \frac{\pi}{2}\right] \) 上的值域.
\end{enumerate}
\end{problem}

\begin{solution}
由积化和差公式，
\[
2\sin\left(x-\frac\pi4\right)\sin\left(x+\frac\pi4\right)
=\cos\left(-\frac\pi2\right)-\cos2x=-\cos2x.
\]
所以
\[
f(x)=\cos\left(2x-\frac\pi3\right)-\cos2x
=\sin\left(2x-\frac\pi6\right).
\]
因此函数的最小正周期为 \(T=\pi\). 图象的对称轴满足
\[
2x-\frac\pi6=\frac\pi2+k\pi\quad(k\in\Z),
\]
即 \(x=\frac\pi3+\frac{k\pi}{2}\).

当 \(x\in\left[-\frac\pi{12},\frac\pi2\right]\) 时，令 \(t=2x-\frac\pi6\)，则 \(t\in\left[-\frac\pi3,\frac{5\pi}6\right]\). 在此区间内，正弦函数在 \(t=\frac\pi2\) 处取得最大值 \(1\)，在左端点取得最小值 \(-\frac{\sqrt3}{2}\)，故值域为 \(\left[-\frac{\sqrt3}{2},1\right]\).
\end{solution}

\begin{problem}
如图，在四棱锥 \(O - ABCD\) 中，底面 \(ABCD\) 四边长为 \(1\) 的菱形，\(\angle ABC = \frac{\pi}{4}\)，\(OA \perp\) 底面 \(ABCD\)，\(OA = 2\)，\(M\) 为 \(OA\) 的中点，\(N\) 为 \(BC\) 的中点.

\begin{enumerate}
\item 证明：直线 \(MN \parallel\) 平面 \(OCD\)；
\item 求异面直线 \(AB\) 与 \(MD\) 所成角的大小；
\item 求点 \(B\) 到平面 \(OCD\) 的距离.
\end{enumerate}

\bitmapfigure[max width=0.34\linewidth]{img/2008/anhui_science/q18_fig1.png}
\end{problem}

\begin{solution}
建立直角坐标系，使 \(A=(0,0,0)\)，底面为 \(z=0\)，取 \(B=(1,0,0)\). 由菱形边长为 \(1\) 且 \(\angle ABC=\frac\pi4\)，可取
\[
C=\left(1-\frac{\sqrt2}{2},\frac{\sqrt2}{2},0\right),\qquad
D=\left(-\frac{\sqrt2}{2},\frac{\sqrt2}{2},0\right).
\]
又因为 \(OA=2\)，所以 \(O=(0,0,2)\)，从而
\[
M=(0,0,1),\qquad
N=\left(1-\frac{\sqrt2}{4},\frac{\sqrt2}{4},0\right).
\]

由平行四边形关系 \(D=A+C-B\)，并结合 \(A=(0,0,0)\)，有
\[
\overrightarrow{MN}=\frac{\overrightarrow{AB}+\overrightarrow{AC}-\overrightarrow{AO}}2
=\overrightarrow{OC}-\frac12\overrightarrow{OD}.
\]
右端是平面 \(OCD\) 内两条相交直线方向向量的线性组合，故 \(MN\parallel\) 平面 \(OCD\).

\(\overrightarrow{AB}=(1,0,0)\)，\(\overrightarrow{MD}=\left(-\frac{\sqrt2}{2},\frac{\sqrt2}{2},-1\right)\)，所以
\[
\left|\overrightarrow{AB}\cdot\overrightarrow{MD}\right|=\frac{\sqrt2}{2},\qquad
|AB|=1,\qquad |MD|=\sqrt2.
\]
设异面直线所成的锐角为 \(\theta\)，则
\[
\cos\theta=\frac{\left|\overrightarrow{AB}\cdot\overrightarrow{MD}\right|}{|AB|\,|MD|}=\frac12,
\]
故 \(\theta=\frac\pi3\).

取 \(\overrightarrow{OC}\) 和 \(\overrightarrow{OD}\) 的向量积，可得平面 \(OCD\) 的一个法向量为 \((0,4,\sqrt2)\)，其方程为
\[
4y+\sqrt2z-2\sqrt2=0.
\]
所以点 \(B=(1,0,0)\) 到该平面的距离为
\[
\frac{|4\cdot0+\sqrt2\cdot0-2\sqrt2|}{\sqrt{4^2+(\sqrt2)^2}}=\frac23.
\]
\end{solution}

\begin{problem}
为防止风沙危害，某地决定建设防护绿化带，种植杨树，沙柳等植物. 某人一次种植了 \(n\) 株沙柳，各株沙柳成活与否是相互独立的，成活率为 \(p\). 设 \(\xi\) 为成活沙柳的株数，数学期望 \(E_{\xi} = 3\)，标准差 \(\sigma_{\xi}\) 为 \(\frac{\sqrt{6}}{2}\).

\begin{enumerate}
\item 求 \(n\)，\(p\) 的值，并写出 \(\xi\) 的分布列；
\item 若有 \(3\) 株或 \(3\) 株以上的沙柳未成活，则需要补种，求需要补种沙柳的概率.
\end{enumerate}
\end{problem}

\begin{solution}
各株沙柳是否成活相互独立，故 \(\xi\sim B(n,p)\). 由二项分布的期望和方差公式，
\[
np=E\xi=3,\qquad np(1-p)=\sigma_{\xi}^{2}=\frac32.
\]
两式相除得 \(1-p=\frac12\)，所以 \(p=\frac12\)，再由 \(np=3\) 得 \(n=6\).

因此对 \(k=0,1,\ldots,6\)，
\[
P(\xi=k)=\mathrm{C}_6^k\left(\frac12\right)^k\left(\frac12\right)^{6-k}=\frac{\mathrm{C}_6^k}{64}.
\]
分布列为
\begin{center}
\begin{tabular}{c|ccccccc}
 \(k\)&\(0\)&\(1\)&\(2\)&\(3\)&\(4\)&\(5\)&\(6\)\\ \hline
 \(P(\xi=k)\)&\(\frac1{64}\)&\(\frac6{64}\)&\(\frac{15}{64}\)&\(\frac{20}{64}\)&\(\frac{15}{64}\)&\(\frac6{64}\)&\(\frac1{64}\)
\end{tabular}
\end{center}

未成活的株数为 \(6-\xi\). 需要补种当且仅当 \(6-\xi\geqslant3\)，即 \(\xi\leqslant3\)，所以所求概率为
\[
P(\xi\leqslant3)=\frac{\mathrm{C}_6^0+\mathrm{C}_6^1+\mathrm{C}_6^2+\mathrm{C}_6^3}{64}
=\frac{1+6+15+20}{64}=\frac{21}{32}.
\]
\end{solution}

\begin{problem}
设函数 \( f(x) = \frac{1}{x \ln x} \)（\(x > 0\) 且 \(x \neq 1\)）.

\begin{enumerate}
\item 求函数 \( f(x) \) 的单调区间；
\item 已知 \(2^{\frac{1}{x}} > x^{a}\) 对任意 \(x \in (0,1)\) 成立，求实数 \(a\) 的取值范围.
\end{enumerate}
\end{problem}

\begin{solution}
函数定义域为 \((0,1)\cup(1,+\infty)\)，且
\[
f'(x)=-\frac{\ln x+1}{(x\ln x)^2}.
\]
因为分母恒为正，所以在 \((0,\e^{-1})\) 上 \(f'(x)>0\)，在 \((\e^{-1},1)\) 和 \((1,+\infty)\) 上 \(f'(x)<0\). 因此 \(f(x)\) 的单调递增区间为 \((0,\e^{-1})\)，单调递减区间为 \((\e^{-1},1)\) 和 \((1,+\infty)\).

对任意 \(x\in(0,1)\)，两边均为正，取自然对数得
\[
\frac{\ln2}{x}>a\ln x.
\]
由于 \(\ln x<0\)，上式等价于
\[
a>\frac{\ln2}{x\ln x}.
\]
由第一问，函数 \(\frac{1}{x\ln x}\) 在 \((0,1)\) 上的最大值为
\[
\frac{1}{\e^{-1}\ln(\e^{-1})}=-\e.
\]
故右端的最大值为 \(-\e\ln2\). 因为原不等式是严格不等式，且在 \(x=\e^{-1}\) 处取得这个最大值，所以必须且只需
\[
a>-\e\ln2.
\]
\end{solution}

\begin{problem}
设数列 \(\{a_{n}\}\) 满足 \(a_0 = 0\)，\(a_{n+1} = c a_n^3 + 1 - c\)，\(n \in \N^*\)，其中 \(c\) 为实数.

\begin{enumerate}
\item 证明：\(a_{n} \in [0,1]\) 对任意 \(n \in \N^*\) 成立的充分必要条件是 \(c \in [0,1]\)；
\item 设 \(0 < c < \frac{1}{3}\). 证明：\(a_{n} \geqslant 1 - (3c)^{n-1}\)，\(n \in \N^*\)；
\item 设 \(0 < c < \frac{1}{3}\). 证明：\(a_1^2 + a_2^2 + \cdots + a_n^2 > n + 1 - \frac{2}{1 - 3c}\)，\(n \in \N^*\).
\end{enumerate}
\end{problem}

\begin{solution}
由递推式取 \(n=0\)，得 \(a_1=1-c\).

若 \(c\in[0,1]\)，则 \(a_1\in[0,1]\). 假设 \(a_n\in[0,1]\)，则
\[
0\leqslant ca_n^3\leqslant c,\qquad 1-c\leqslant a_{n+1}=ca_n^3+1-c\leqslant1,
\]
所以 \(a_{n+1}\in[0,1]\). 由数学归纳法，\(a_n\in[0,1]\) 对任意 \(n\in\N^*\) 成立.

反之，若对任意 \(n\in\N^*\) 都有 \(a_n\in[0,1]\)，则 \(a_1=1-c\in[0,1]\)，从而 \(c\in[0,1]\). 因此所述条件的充分必要条件是 \(c\in[0,1]\).

以下设 \(0<c<\frac13\). 由第一问知 \(0\leqslant a_n\leqslant1\). 先证第二问：当 \(n=1\) 时，\(a_1=1-c\geqslant0=1-(3c)^0\). 若 \(a_n\geqslant1-(3c)^{n-1}\)，则
\[
1-a_{n+1}=c(1-a_n^3)=c(1-a_n)(1+a_n+a_n^2)
\leqslant3c(1-a_n)\leqslant(3c)^n.
\]
故 \(a_{n+1}\geqslant1-(3c)^n\)，结论成立.

令 \(q=3c\in(0,1)\). 由第二问，且 \(a_k\geqslant0\)，有
\[
a_k^2\geqslant\left(1-q^{k-1}\right)^2
=1-2q^{k-1}+q^{2k-2}.
\]
于是
\[
\begin{aligned}
\sum_{k=1}^{n}a_k^2
&\geqslant n-\frac{2(1-q^n)}{1-q}+\frac{1-q^{2n}}{1-q^2}\\
&=n+1-\frac{2}{1-q}
+\frac{2q^n}{1-q}+\frac{1-q^{2n}}{1-q^2}-1.
\end{aligned}
\]
其中 \(\frac{2q^n}{1-q}>0\)，且 \(\frac{1-q^{2n}}{1-q^2}\geqslant1\)，故
\[
\sum_{k=1}^{n}a_k^2>n+1-\frac{2}{1-q}=n+1-\frac{2}{1-3c}.
\]
\end{solution}

\begin{problem}
设椭圆 \(C: \frac{x^2}{a^2} + \frac{y^2}{b^2} = 1\)（\(a > b > 0\)）过点 \(M(\sqrt{2}, 1)\). 且左焦点为 \(F_1(-\sqrt{2}, 0)\).

\begin{enumerate}
\item 求椭圆 \(C\) 的方程；
\item 当过点 \(P(4,1)\) 的动直线 \(l\) 与椭圆 \(C\) 相交于两不同点 \(A\)，\(B\) 时，在线段 \(AB\) 上取点 \(Q\)，满足 \(\left|\overrightarrow{AP}\right| \cdot \left|\overrightarrow{QB}\right| = \left|\overrightarrow{AQ}\right| \cdot \left|\overrightarrow{PB}\right|\)，证明：点 \(Q\) 总在某定直线上.
\end{enumerate}
\end{problem}

\begin{solution}
\begin{enumerate}
\item 设椭圆的半焦距为 \(c\). 由左焦点为 \(F_1(-\sqrt2,0)\)，得 \(c=\sqrt2\)，所以
\[
a^2-b^2=c^2=2.
\]
又因为点 \(M(\sqrt2,1)\) 在椭圆上，故
\[
\frac{2}{a^2}+\frac{1}{b^2}=1.
\]
令 \(A=a^2\)，则 \(b^2=A-2\)，代入上式得
\[
\frac2A+\frac1{A-2}=1
\quad\Longrightarrow\quad
2(A-2)+A=A(A-2),
\]
即 \(A^2-5A+4=0\). 因而 \(A=1\) 或 \(A=4\). 由 \(a>b>0\) 且 \(a^2-b^2=2\)，只能取 \(a^2=4\)，于是 \(b^2=2\). 椭圆方程为
\[
\frac{x^2}{4}+\frac{y^2}{2}=1.
\]

\item 设过 \(P(4,1)\) 的动直线的方向向量为 \(\bs{u}=(\alpha,\beta)\)，用参数 \(t\) 表示直线上的点：
\[
X=P+t\bs{u}.
\]
设直线与椭圆的交点对应的参数为 \(t_1,t_2\)，并取标号使 \(0<t_1<t_2\). 将 \(x=4+t\alpha\)、\(y=1+t\beta\) 代入椭圆方程，得
\[
\left(\frac{\alpha^2}{4}+\frac{\beta^2}{2}\right)t^2+(2\alpha+\beta)t+\frac72=0.
\]
记 \(A_0=\frac{\alpha^2}{4}+\frac{\beta^2}{2}\)，由韦达定理，
\[
t_1+t_2=-\frac{2\alpha+\beta}{A_0},\qquad
t_1t_2=\frac{7}{2A_0}.
\]

点 \(Q\) 在线段 \(AB\) 上，题设条件等价于
\[
\frac{|AQ|}{|QB|}=\frac{|AP|}{|PB|}=\frac{t_1}{t_2}.
\]
若 \(Q=P+t_Q\bs{u}\)，则 \(|AQ|=(t_Q-t_1)|\bs{u}|\)、\(|QB|=(t_2-t_Q)|\bs{u}|\)，所以
\[
t_1(t_2-t_Q)=t_2(t_Q-t_1),\qquad
t_Q=\frac{2t_1t_2}{t_1+t_2}=-\frac7{2\alpha+\beta}.
\]
这里 \(2\alpha+\beta\neq0\)，否则上面的交点方程无两个不同实根. 因而若 \(Q=(x,y)\)，则
\[
x=4-\frac{7\alpha}{2\alpha+\beta},\qquad
y=1-\frac{7\beta}{2\alpha+\beta}.
\]
于是
\[
2x+y-2=7-\frac{7(2\alpha+\beta)}{2\alpha+\beta}=0.
\]
所以点 \(Q\) 总在定直线 \(2x+y-2=0\) 上.
\end{enumerate}
\end{solution}
